% from ega0-1-fr.tex
\oldpage[0\textsubscript{I}]{63}

d'idéaux. Les applications canoniques
$\varphi_\lambda:A\to A/\mathfrak J_\lambda$ forment un système projectif de
représentations continues et définissent donc une représentation continue
$\varphi:A\to\varprojlim A/\mathfrak J_\lambda$~; si $A$ est séparé,
$\varphi$ est un isomorphisme topologique de $A$ sur un
sous-anneau partout dense de $\varprojlim A/\mathfrak J_\lambda$~; si, en
outre, $A$ est complet, $\varphi$ est un isomorphisme topologique de $A$ sur
$\varprojlim A/\mathfrak J_\lambda$.
\end{env}

\begin{lemma}[7.2.2]
\label{0.7.2.2-fr}
Pour qu'un anneau linéairement topologisé soit admissible, il faut et il suffit
qu'il soit isomorphe à une limite projective
$A=\varprojlim A_\lambda$, où $(A_\lambda,u_{\lambda\mu})$ est un système
projectif d'anneaux discrets ayant pour ensemble d'indices un ensemble ordonné
filtrant $L$ (pour $\leq$) qui admet un plus petit élément noté $0$ et satisfait
aux conditions suivantes~: 1\textsuperscript{o} les
$u_\lambda:A\to A_\lambda$ sont surjectifs~; 2\textsuperscript{o} le noyau
$\mathfrak J_\lambda$ de $u_{0\lambda}:A_\lambda\to A_0$ est nilpotent.
Lorsqu'il en est ainsi, le noyau $\mathfrak J$ de $u_0:A\to A_0$ est égal à
$\varprojlim\mathfrak J_\lambda$.
\end{lemma}

La nécessité de la condition résulte de (7.2.1), en prenant pour
$(\mathfrak J_\lambda)$ un système fondamental de voisinages de $0$ formé
d'idéaux de définition contenus dans l'un d'eux $\mathfrak J_0$ et appliquant
(7.1.4, (i)). La réciproque résulte de la définition d'une limite projective et
de (7.1.2), et la dernière assertion est immédiate.

\begin{env}[7.2.3]
\label{0.7.2.3-fr}
Soient $A$ un anneau topologique admissible, $\mathfrak J$ un idéal de $A$
contenu dans un idéal de définition (autrement dit (7.1.4) tel que
$(\mathfrak J^n)$ tende vers $0$)~; on peut considérer sur $A$ la topologie
d'anneau ayant pour système fondamental de voisinages de $0$ les puissances
$\mathfrak J^n$ ($n>0$)~; nous l'appellerons encore la topologie
\emph{$\mathfrak J$-préadique}. L'hypothèse que $A$ est admissible entraîne que
$\bigcap_n\mathfrak J^n=0$, donc la topologie $\mathfrak J$-préadique sur $A$
est séparée~; soit
$\widehat A=\varprojlim A/\mathfrak J^n$ le complété de $A$ pour cette topologie
(où les $A/\mathfrak J^n$ sont munis de la topologie discrète), et désignons
par $u$ l'homomorphisme d'anneaux $A\to\widehat A$ (non nécessairement
continu), limite projective de la suite d'homomorphismes
$u_n:A\to A/\mathfrak J^n$. D'autre part, la topologie $\mathfrak J$-préadique
sur $A$ est plus fine que la topologie donnée $\mathcal T$ sur $A$~; comme $A$
est séparé et complet pour $\mathcal T$, on peut prolonger par continuité
l'application identique de $A$ (muni de la topologie $\mathfrak J$-préadique)
dans $A$ muni de $\mathcal T$~; cela donne une représentation continue
$v:\widehat A\to A$.
\end{env}

\begin{proposition}[7.2.4]
\label{0.7.2.4-fr}
Si $A$ est un anneau admissible et $\mathfrak J$ est contenu dans un idéal de
définition de $A$, $A$ est séparé et complet pour la topologie
$\mathfrak J$-préadique.
\end{proposition}

En effet, avec les notations de (7.2.3), il est immédiat que $v\circ u$ est
l'application identique de $A$. D'autre part,
$u_n\circ v:\widehat A\to A/\mathfrak J^n$ est le prolongement par continuité
(pour la topologie $\mathfrak J$-préadique sur $A$ et la topologie discrète sur
$A/\mathfrak J^n$) de l'application canonique $u_n$~; autrement dit, c'est
l'application canonique de
$\widehat A=\varprojlim_k A/\mathfrak J^k$ sur $A/\mathfrak J^n$~;
$u\circ v$ est donc limite projective de cette suite d'applications,
c'est-à-dire par définition l'application identique de $\widehat A$~; ceci
démontre la proposition.

\begin{corollary}[7.2.5]
\label{0.7.2.5-fr}
Sous les hypothèses de (7.2.3), les conditions suivantes sont équivalentes~:
\begin{enumerate}
  \item[$a)$] L'homomorphisme $u$ est continu~;

